\section{Legal Landscape}
\label{sec:legal}

\subsection{The Symmetry Standard in Academic Literature}

The symmetry standard was developed by \citet{gelmanKing1994} as a normative
framework for evaluating partisan fairness in redistricting. Their core argument
is that equal treatment of political parties requires symmetric treatment in the
seats-votes function: if a party receives $v\%$ of votes, it should win the same
seat share that any other party would win at $v\%$ of votes. This symmetry principle
is grounded in equal protection and democratic theory, not in proportionality.

Partisan Bias, as the scalar summary of symmetry at the 50\% pivot, is the most
commonly reported element of the symmetry standard in litigation. \citet{gelmanKing1994}
argue that courts can operationalize the symmetry standard by comparing the seats-votes
curve to the anti-symmetric ideal and reporting Partisan Bias as the summary statistic.
Unlike the Efficiency Gap, the symmetry standard does not propose a specific threshold;
it measures a directional property (partisan advantage at the pivot) whose magnitude
is relevant to the constitutional analysis.

\subsection{Federal Court History}

Partisan Bias under the symmetry standard was cited in North Carolina and Wisconsin
redistricting litigation in the years between \textit{Gill v.\ Whitford} (2018)
and \textit{Rucho v.\ Common Cause} (2019). In the \textit{Common Cause v.\ Rucho}
district court proceedings,\footnote{\textit{Common Cause v.\ Rucho}, 318 F.Supp.3d
777 (M.D.N.C.\ 2018).} expert witnesses reported Partisan Bias as part of a
multi-metric analysis. The district court credited the symmetry standard evidence
and ruled for the plaintiffs; SCOTUS reversed in \textit{Rucho} on political
question grounds without addressing the merits of the symmetry standard.

\textit{Rucho}'s majority opinion noted the symmetry standard as one of several
proposed metrics that fail to provide a judicially manageable standard, primarily
on the grounds that the Court could not identify what level of asymmetry constitutes
a constitutional violation. This critique is methodological (no threshold) rather
than empirical (the measurement itself is rejected). Expert witnesses in state
court proceedings can continue to use Partisan Bias and the symmetry standard;
the threshold question is a matter of state constitutional interpretation that
\textit{Rucho} explicitly left to state courts.

\subsection{State Court Applications}

\textbf{North Carolina} (\textit{Harper v.\ Hall}, 2022): The NC Supreme Court
cited symmetry evidence, including Partisan Bias, in its finding that the enacted
congressional map systematically advantaged Republicans. Enacted NC Partisan Bias
of approximately $-0.07$ (one-seat Republican advantage) was part of the evidentiary
record that the court credited.

\textbf{Pennsylvania}: PA redistricting proceedings have included expert testimony
on the symmetry standard. The PA Supreme Court's analysis in \textit{League of
Women Voters} focused more on the comparative analysis of proposed maps than on
specific metric thresholds, but the symmetry principle informed the court's reasoning.

\textbf{Ohio}: Ohio redistricting proceedings included symmetry evidence in support
of claims that the Ohio congressional map violated the state constitution's explicit
anti-gerrymandering provision.

\subsection{The Jury-Friendly Argument}

Partisan Bias's strongest courtroom advantage is its communicative simplicity.
The argument is:
\begin{quote}
\textit{Under the enacted NC map, if both parties received exactly 50\% of
the two-party vote, Republicans would win 8 seats and Democrats would win 6 seats.
Under the bisect algorithmic map, both parties would win 7 seats at 50\% of the vote.
The difference---one additional Republican seat guaranteed by the map geometry, not
by voter choices---is the Partisan Bias.}
\end{quote}
This argument does not require the court to understand wasted-vote algebra, mean-median
statistics, or swing model assumptions. It communicates partisan asymmetry directly in
terms of seat outcomes at equal vote shares, which is precisely the equal-treatment
principle that courts applying ``free and equal elections'' standards are designed to enforce.
